\documentclass[12pt,a4paper]{article}

\usepackage[T2A]{fontenc}
\usepackage[utf8]{inputenc}
\usepackage[russian]{babel}

\usepackage{graphicx}
\usepackage{amsmath, amssymb}
\usepackage{geometry}
\geometry{margin=2.3cm}

\title{1D HJB + nD LQR}
\date{May 2026}


\begin{document}

\maketitle

\section{Постановка задачи}

Пусть $c_t=\mathrm{CAR}_t$. Считаем, что из одномерной HJB-задачи уже получена внешняя политика общего роста портфеля
\[
    u(c_t).
\]
В nD-задаче это не управление, а заданная функция. Управляем только структурой портфеля.

Общий размер портфеля обозначим через
\[
    Q_t = \sum_i q_t^i.
\]

Динамика общего размера портфеля (шум не убираем):
\[
    dQ_t = Q_t\left(u(c_t)dt + \nu_Q dW_t^Q\right).
\]

Доли активов в портфеле:
\[
    r_t^i = \frac{q_t^i}{Q_t},
    \qquad
    \sum_i r_t^i = 1.
\]

Динамика долей:
\[
    dr_t^i = r_t^i \theta_t^i dt.
\]

Так как сумма долей должна сохраняться, имеем ограничение:
\[
    \sum_i r_t^i \theta_t^i = 0.
\]

Поскольку
\[
    q_t^i = Q_t r_t^i,
\]
то управляемая часть относительного изменения позиции $q_t^i$ имеет вид
\[
    \frac{dq_t^i}{q_t^i}
    =
    \frac{dQ_t}{Q_t}
    +
    \frac{dr_t^i}{r_t^i}
    =
    \left(u(c_t)+\theta_t^i\right)dt
    + \nu_Q dW_t^Q.
\]

Liquidity cost считаем только по управляемой части скорости изменения позиции:
\[
    u(c_t)+\theta_t^i.
\]

Тогда динамика прибыли портфеля:
\[
\begin{aligned}
    dG_t
    =
    &\; Q_t r_t^\top \mu \, dt \\
    &- Q_t \sum_i \lambda_i r_t^i
        \left(u(c_t)+\theta_t^i\right)^2 dt \\
    &+ Q_t (r_t \odot \sigma_t)^\top dP_t.
\end{aligned}
\]

Раскроем liquidity-слагаемое:
\[
\begin{aligned}
    Q_t \sum_i \lambda_i r_t^i
        \left(u(c_t)+\theta_t^i\right)^2
    =
    &\; Q_t u(c_t)^2 \sum_i \lambda_i r_t^i \\
    &+ 2Q_t u(c_t) \sum_i \lambda_i r_t^i\theta_t^i \\
    &+ Q_t \sum_i \lambda_i r_t^i(\theta_t^i)^2.
\end{aligned}
\]

То есть динамику прибыли (валовая прибыль) можно записать как
\[
\begin{aligned}
    dG_t
    =
    &\; Q_t r_t^\top \mu \, dt \\
    &- Q_t u(c_t)^2 \lambda^\top r_t \, dt \\
    &- 2Q_t u(c_t)\lambda^\top(r_t \odot \theta_t)\,dt \\
    &- Q_t\sum_i \lambda_i r_t^i(\theta_t^i)^2dt \\
    &+ Q_t (r_t \odot \sigma_t)^\top dP_t.
\end{aligned}
\]

Здесь $\lambda=(\lambda_1,\dots,\lambda_n)$, а $\odot$ обозначает покомпонентное произведение.

Первое liquidity-слагаемое отвечает за стоимость общего роста или сокращения всего портфеля. Второе слагаемое является cross-term между общим ростом портфеля и изменением структуры. Третье слагаемое отвечает за стоимость изменения структуры портфеля.

Капитал:
\[
    dK_t = (1-h)dG_t.
\]

Дивиденды:
\[
    dD_t = h dG_t.
\]

Здесь $h$ -- доля прибыли, которая выплачивается в виде дивидендов.

Лагранжиан зададим как поток дивидендов (максимизируем заработок акционера) с риск-штрафом:
\[
    dL_t
    =
    h dG_t
    -
    \frac{1}{2}\eta Q_t r_t^\top \Sigma r_t \, dt.
\]

Функционал:
\[
    V(t, state)
    =
    \sup_{\theta}
    \mathbb{E}
    \left[
        \int_t^T e^{-\rho(s-t)} dL_s
        +
        g(state_T)
    \right].
\]

Ограничения:
\[
    \sum_i r_t^i = 1,
    \qquad
    r_t^i \ge 0,
\]
\[
    \sum_i r_t^i\theta_t^i = 0.
\]

Также при необходимости можно добавить box constraints:
\[
    \theta_{\min}^i \le \theta_t^i \le \theta_{\max}^i.
\]

\subsection{Стейт в задаче}

В данной постановке состояние можно записать как
\[
    X_t = (Q_t, K_t, r_t),
\]
где $Q_t$ -- общий размер портфеля, $K_t$ -- капитал банка, а
\[
    r_t = (r_t^1,\dots,r_t^n)
\]
-- вектор долей активов в портфеле. Текущий уровень достаточности капитала задается как
\[
    c_t = \mathrm{CAR}_t.
\]
Если $\mathrm{CAR}$ выражается через капитал и общий размер портфеля, например
\[
    c_t = \frac{K_t}{\omega Q_t},
\]
то вместо состояния $(Q_t,K_t,r_t)$ можно использовать эквивалентное состояние
\[
    \widetilde X_t = (Q_t,c_t,r_t).
\]
При этом вектор долей лежит на симплексе:
\[
    \sum_i r_t^i = 1,
    \qquad
    r_t^i \ge 0.
\]

\subsection{Сведение к LQR}

В текущих переменных задача не является LQR-задачей. Динамика долей имеет вид
\[
    dr_t^i = r_t^i\theta_t^i dt,
\]
то есть содержит произведение состояния $r_t^i$ и управления $\theta_t^i$. В классической LQR-задаче динамика должна быть линейной по состоянию и управлению.

Чтобы получить линейную динамику, введем новую переменную управления
\[
    v_t^i = r_t^i\theta_t^i.
\]
Тогда
\[
    dr_t^i = v_t^i dt,
\]
или в векторной форме
\[
    \dot r_t = v_t.
\]
Ограничение на сохранение суммы долей также становится линейным:
\[
    \sum_i r_t^i\theta_t^i = 0
    \quad \Longleftrightarrow \quad
    \sum_i v_t^i = 0.
\]

Однако после замены в функционале возникает нелинейное слагаемое. Действительно (на длинном горизонте берем стационарную точку),
\[
    Q_t\sum_i \lambda_i r_t^i(\theta_t^i)^2
    =
    Q_t\sum_i \lambda_i r_t^i
    \left(\frac{v_t^i}{r_t^i}\right)^2
    =
    Q_t\sum_i \lambda_i\frac{(v_t^i)^2}{r_t^i}.
\]
Это выражение не является квадратичным по паре переменных $(r_t^i,v_t^i)$ из-за деления на $r_t^i$.

Чтобы получить LQR/QP-аппроксимацию, разложим это слагаемое локально в окрестности фиксированной точки $r^{\mathrm{ref},0}$. Для этого заморозим знаменатель:
\[
    \frac{1}{r_t^i}
    \approx
    \frac{1}{r_i^{\mathrm{ref},0}}.
\]
Тогда
\[
    Q_t\sum_i \lambda_i\frac{(v_t^i)^2}{r_t^i}
    \approx
    Q_t\sum_i
    \frac{\lambda_i}{r_i^{\mathrm{ref},0}}
    (v_t^i)^2.
\]
Обозначим
\[
    \overline{\Lambda}
    =
    \operatorname{diag}
    \left(
        \frac{\lambda_1}{r_1^{\mathrm{ref},0}},
        \dots,
        \frac{\lambda_n}{r_n^{\mathrm{ref},0}}
    \right).
\]
Тогда приближенное liquidity-слагаемое принимает квадратичный вид
\[
    Q_t v_t^\top \overline{\Lambda} v_t.
\]

Таким образом, после замены $v_t=r_t\odot\theta_t$ и локальной аппроксимации около $r^{\mathrm{ref},0}$ задача по долям получает линейную динамику и квадратичный функционал. Поэтому ее можно рассматривать как constrained LQR, а после дискретизации -- как QP-задачу.

\subsection{QP-задача}

После замены
\[
    v_t^i = r_t^i\theta_t^i
\]
и локальной аппроксимации
\[
    \frac{1}{r_t^i} \approx \frac{1}{r_i^{\mathrm{ref},0}},
\]
получаем дискретную QP-задачу. Пусть временная сетка имеет вид
\[
    t_k = k\Delta t,\qquad k=0,\dots,N.
\]

Динамика долей:
\[
    r_{k+1} = r_k + \Delta t\, v_k.
\]

Ограничения:
\[
    \mathbf{1}^\top r_k = 1,
    \qquad
    r_k^i \ge 0,
\]
\[
    \mathbf{1}^\top v_k = 0.
\]

Обозначим
\[
    \overline{\Lambda}
    =
    \operatorname{diag}
    \left(
        \frac{\lambda_1}{r_1^{\mathrm{ref},0}},
        \dots,
        \frac{\lambda_n}{r_n^{\mathrm{ref},0}}
    \right).
\]

При фиксированном значении внешней политики общего роста портфеля
\[
    \bar u = u(c_t)
\]
QP-задача имеет вид
\[
\begin{aligned}
    \min_{\{r_k,v_k\}_{k=0}^{N-1}}
    \sum_{k=0}^{N-1}
    e^{-\rho t_k}Q_k\Delta t
    \Big[
        &-h\mu^\top r_k
        + h\bar u^2\lambda^\top r_k \\
        &+ 2h\bar u\,\lambda^\top v_k
        + h v_k^\top\overline{\Lambda}v_k \\
        &+ \frac{1}{2}\eta r_k^\top\Sigma r_k
    \Big].
\end{aligned}
\]

Здесь первое слагаемое отвечает за доходность портфеля, второе и третье появляются из раскрытия liquidity cost
\[
    \sum_i \lambda_i r_i(\bar u+\theta_i)^2,
\]
четвертое является квадратичной аппроксимацией стоимости изменения структуры портфеля, а последнее слагаемое задает риск-штраф с риск-аппетитом $\eta$.

После дискретизации задача имеет квадратичную целевую функцию и линейные ограничения, поэтому является QP-задачей.

\subsection{Полный подход}

Итоговый алгоритм состоит из двух частей: offline-этапа и online/simulation-этапа.

\subsubsection*{Идеи}

\begin{enumerate}
    \item Решать 1D HJB онлайн.

    \item Делать риск-аппетит функцией времени или состояния:
    \[
        \eta_t = \eta(t)
        \quad \text{или} \quad
        \eta_t = \eta(c_t,r_t).
    \]
    Отдельный вопрос -- как именно связывать $\eta_t$ с текущим уровнем CAR: PD (PB), через mean-reversion CAR к целевому уровню или через другую метрику риска.
\end{enumerate}

\subsubsection*{Offline-этап}

На offline-этапе строим приближение для политики общего роста портфеля. Для этого берем набор фиксированных структур портфеля
\[
    r_{\mathrm{fixed}} \in r_{\mathrm{grid}}.
\]

Для каждой фиксированной структуры портфеля агрегируем параметры:
\[
    \mu_R = r_{\mathrm{fixed}}^\top \mu,
\]
\[
    \sigma_R^2 = r_{\mathrm{fixed}}^\top \Sigma r_{\mathrm{fixed}},
\]
\[
    \lambda_R = \lambda^\top r_{\mathrm{fixed}},
\]
\[
    w_R = w^\top r_{\mathrm{fixed}}.
\]

То есть вместо полной 20-мерной структуры $r_{\mathrm{fixed}}$ одномерная HJB-задача видит только агрегированный вектор параметров:
\[
    z_R =
    (\mu_R,\sigma_R,\lambda_R,w_R).
\]

Далее для каждого такого набора агрегированных параметров решаем одномерную HJB-задачу по состоянию
\[
    c = \mathrm{CAR}.
\]

На выходе получаем политику общего роста портфеля:
\[
    u(c; z_R).
\]

Таким образом, offline можно построить таблицу вида:
\[
    (c,z_R) \mapsto u(c;z_R).
\]

Можно интерпретировать ее так: по столбцам идут значения CAR, по строкам -- агрегированные параметры портфеля $z_R$, а в ячейках лежит оптимальное значение общего роста портфеля $u$.

Если таблица посчитана не во всех точках, то можно использовать интерполяцию или обучить приближенную модель:
\[
    \widehat u(c,\mu_R,\sigma_R,\lambda_R,w_R).
\]

\subsubsection*{Online / simulation-этап}

Пусть задано начальное состояние:
\[
    Q_0,\qquad K_0,\qquad r_0.
\]

Далее для каждого шага времени $t=0,1,\dots,T-1$ выполняем следующий алгоритм.

\begin{enumerate}
    \item Считаем текущий уровень достаточности капитала:
    \[
        c_t =
        \frac{K_t}{(w^\top r_t)Q_t}.
    \]

    \item Агрегируем текущую структуру портфеля $r_t$:
    \[
        \mu_R(t) = r_t^\top \mu,
    \]
    \[
        \sigma_R^2(t) = r_t^\top \Sigma r_t,
    \]
    \[
        \lambda_R(t) = \lambda^\top r_t,
    \]
    \[
        w_R(t) = w^\top r_t.
    \]

    Обозначим
    \[
        z_t =
        (\mu_R(t),\sigma_R(t),\lambda_R(t),w_R(t)).
    \]

    \item По HJB-таблице или приближенной модели получаем управление общим размером портфеля:
    \[
        \bar u_t =
        \widehat u(c_t,z_t).
    \]

    \item На горизонте LQR/QP считаем $\bar u_t$ фиксированным, то есть замораживаем политику общего роста:
    \[
        u_s \approx \bar u_t,
        \qquad s \in [t,t+T_{\mathrm{LQR}}].
    \]

    Это нужно, чтобы задача по структуре портфеля оставалась QP-задачей.

    \item При фиксированном $\bar u_t$ решаем LQR/QP-задачу по структуре портфеля для сетки риск-аппетитов:
    \[
        \eta \in \eta_{\mathrm{grid}}.
    \]

    Для каждого $\eta$ получаем оптимальную траекторию управлений:
    \[
        v_0^\eta, v_1^\eta,\dots
    \]

    \item Выбираем риск-аппетит $\eta_t$.

    На первом этапе можно задавать его вручную, например:
    \[
        \eta_t = 1.
    \]

    В дальнейшем $\eta_t$ можно выбирать из условия согласования вероятности пробития:
    \[
        PD(\eta_t,r_t) = PD_{\mathrm{target}}(c_t).
    \]

    \item Берем первый шаг LQR-управления:
    \[
        v_t = v_0^{\eta_t}.
    \]

    Это MPC-style подход: задача решается на горизонте, но применяется только первое управление.

    \item Обновляем общий размер портфеля:
    \[
        Q_{t+\Delta t}
        =
        Q_t + Q_t\bar u_t\Delta t + \mathrm{noise}_Q.
    \]

    Если шум в $Q_t$ не учитывается, используем детерминированное обновление:
    \[
        Q_{t+\Delta t}
        =
        Q_t\exp(\bar u_t\Delta t).
    \]

    \item Обновляем структуру портфеля:
    \[
        r_{t+\Delta t}
        =
        r_t + v_t\Delta t.
    \]

    При этом должно сохраняться условие:
    \[
        \sum_i r_{t+\Delta t}^i = 1.
    \]

    \item Считаем прибыль портфеля:
    \[
    \begin{aligned}
        dG_t
        =
        &\; Q_t r_t^\top \mu \, \Delta t \\
        &- Q_t\sum_i \lambda_i r_t^i
        \left(\bar u_t+\theta_t^i\right)^2\Delta t \\
        &+ \mathrm{market\ noise}.
    \end{aligned}
    \]

    Здесь
    \[
        \theta_t^i = \frac{v_t^i}{r_t^i}.
    \]

    \item Обновляем капитал:
    \[
        K_{t+\Delta t}
        =
        K_t + (1-h)dG_t.
    \]

    \item Обновляем накопленные дивиденды:
    \[
        D_{t+\Delta t}
        =
        D_t + h dG_t.
    \]

    \item Считаем новый CAR:
    \[
        c_{t+\Delta t}
        =
        \frac{K_{t+\Delta t}}
        {(w^\top r_{t+\Delta t})Q_{t+\Delta t}}.
    \]

    \item Проверяем пробитие:
    \[
        \text{if } c_{t+\Delta t} \le c_{\mathrm{bar}},
        \text{ then stop path.}
    \]
\end{enumerate}

Итого: HJB отвечает за общий рост или сокращение портфеля, а LQR/QP отвечает за изменение структуры портфеля. На каждом online-шаге HJB дает одно число $\bar u_t$, затем это значение замораживается, после чего решается LQR/QP-задача по долям.

\subsection{Быстрый constrained LQR}
Вспомним $QP$ задачу. Динамика:
$$
    dr_t = v_t dt, \sum_i r_t^i = 1, r_t^i \ge 0, \sum_i v_t^i = 0
$$
Лагранжиан 

$$L(t, r_t, v_t) = Q_t \cdot 
\left( r_t^\top \mu - 0.5 \eta r_t^\top \Sigma r_t - v_t^\top \overline{\Lambda} v \right)$$
Здесь считаем $Q_t$ -- заданная функция времени. Если считать, что $u_R = const$ -- управление суммарным объёмом постоянное, то
$Q_t = Q_0 e^{u_R t}$. 

Value-function:
$$
    V(t, r_t) = \sup_{v} \int_t^\infty e^{-\rho(s-t)} L(s, r_s, v_s) ds
    = Q_0 \sup_{v} \int_t^\infty e^{-(\rho-u_R)(s-t)} \left( r_s^\top \mu - 0.5 \eta r_s^\top \Sigma r_s - v_s^\top \overline{\Lambda} v_s \right) ds
$$

Т.е. эффективно решаем задачу LQR с ограничением на управление: $e^\top u = 0$:
$$
\begin{aligned}
    &dr_t = v_t dt, e^\top v_t = 0 \\ 
    &\tilde{L}(r_t, v_t) = r_t^\top \mu - 0.5 \eta r_t^\top \Sigma r_t - v_t^\top \overline{\Lambda} v_t \\
\end{aligned}
$$
Вэлью-функция:
$$
    V(x) = \sup_u \int_0^\infty e^{-\tilde{\rho} t} \tilde{L}(r_t, v_t) dt
$$
где $\tilde{\rho} = \rho - u_R$ -- эффективный дисконт-фактор.
Ищем решение в виде $V(x) = r_t^T P r_t + p^T r_t + c$, получаем уравнение Риккати на $P$ и $p$.
$$
    \tilde{\rho} P = P M P - 0.5  \eta \Sigma
$$где $M = \tilde{L}^{-1} - \dfrac{\tilde{L}^{-1}ee^\top \tilde{L}^{-1}}{e\top \tilde{L}^{-1}e}$.
Оптимальная политика:
$$
    v_t = MP (r_t- r^*)
$$где $r^*$ -- оптимальный по Марковицу портфель.

В исходной задаче у нас ограничение на состояние $r$ : $r^i \in [x_{min}, x_{max}]$. 
В таком случае задача на вэлью не сводится к алгебраическому решению риккати -> необходимо решать задачу QP для каждой начальной точки $x_0$ отдельно.

Анзатц: положим оптимальное управление $v_t = MP (r_t - \overline{r})$,
где $M, P$ -- матрицы из решения задачи LQR без ограничения на стейт, $\overline{r}$ -- оптимальный по марковицу портфель при условии ограничений $r^i \in [r_{min}, r_{max}]$.
Получим некоторое прокси оптимального решения. Заметим, что оно совпадает с задачей QP, если ограничения non-binding на оптимальной траектории (в этом случае задача оптимизации с и без ограничений совпадают).

$\overline{r}$ зависит только от $\mu, \eta, \Sigma$. В задаче у нас $\mu, \Sigma$ фиксированны, $\eta$ потенциально может меняться.
Марковица можем в офлайне предпосчитать для сетки по $\eta$, дальше интерполировать.

Матрицы $M, P$ можем предпосчитать в офлайне (если считаем, что косты за ликвидность не зависят от текущей структуры), либо решать в онлайне, что в любом случае намного быстрее, чем решать QP.





\end{document}